\section{Results}

The regenerated benchmark now meaningfully differentiates the defenses. No Protection blocks 0.0\% of the security-deny scenarios. Static Allowlist reaches 50.6\%, Progent-style rules reach 67.47\%, Validator Only reaches 72.29\%, and Argument Allowlist reaches 83.13\%. Full Agent-Sentinel reaches 100.0\%, and the No Audit ablation also reaches 100.0\% because it preserves the same enforcement path while disabling observability rather than authorization.

The category-level breakdown explains these gaps. Static allowlists deny the POST-based policy-blocked family, but they do not constrain private reads, path-escape writes, or disallowed-host GET requests. Validator-only enforcement blocks private reads and disallowed-host GETs, but it does not enforce the stronger policy-only checks that separate full mediation from pure validation. The Progent-style baseline improves over static allowlists through rule-based argument checks, but it still admits some requests that Agent-Sentinel denies through capability policy evaluation. Argument Allowlist performs best among the weaker baselines because it combines tool and argument constraints, but it still falls short of full capability mediation.

Latency remains small in absolute terms. Agent-Sentinel records 0.27\,ms median decision latency and 0.67\,ms p95 latency on the benchmark, while the weaker baselines are faster because they evaluate fewer checks. The key paper result is therefore not that all mediated systems are identical, but that richer runtime mediation produces a clear increase in blocking effectiveness while remaining operationally lightweight.
